\section*{Orientaciones bibliográficas}
\addcontentsline{toc}{section}{Orientaciones bibliográficas}
Para grupos, órdenes y productos, véase \citet[cap.~II, definición 1.2, resultados 1.9--10 y 1.12--14]{aluffi2009}.

Para acciones y grupos libres, véase \citet[cap.~II, 4.1/3/5--8 y 5.1/3]{aluffi2009}. La reducción de palabras y la propiedad universal se desarrollan también en \citet[cap.~2, proposiciones 2.1--3, pp.~32--33]{milne2025}.

Para grupos abelianos libres, subgrupos, núcleos y subgrupos cíclicos, véase \citet[cap.~II, 6.1--4, 6.6, 6.9 y 6.11]{aluffi2009}.

Para cocientes e isomorfismos, véase \citet[cap.~II, 7.1--4/6--10/12 y 8.1--2]{aluffi2009}. El índice y el recuento de clases laterales preparan el estudio de acciones de la sección~\ref{sec:grupos-acciones}.

Para acciones, conjugación y ecuación de clases, véase \citet[cap.~II, 9.6--12, y cap.~IV, 1.1/3/5--7/9/11--14]{aluffi2009}. Cauchy es el teorema IV.2.1; el ejercicio sobre $G/Z(G)$ corresponde a IV.1.5.

Para los teoremas de Sylow y sus aplicaciones, véase \citet[cap.~IV, 2.2--3/5--6/8--16]{aluffi2009}. La comprobación de puntos fijos es el ejercicio IV.1.1.

Para Jordan--Hölder, conmutadores y resolubilidad, véase \citet[cap.~IV, 3.1--13]{aluffi2009}.

Para ciclos, signo y grupos alternantes, véase \citet[cap.~IV, 4.1--4/6/8--21]{aluffi2009}. La unicidad de la descomposición en ciclos es el ejercicio IV.4.2.

Para extensiones, productos semidirectos y clasificación abeliana, véase \citet[cap.~IV, 5.1--3/5--13, 6.1--6 y ejemplo 6.8]{aluffi2009}. El ejemplo del producto semidirecto corresponde al ejercicio IV.5.13.
